\section{Completed Native30 bicoherence extension}
\label{sec:native30-bc-completed}
The completed extension evaluates all 255 nonempty subsets of eight families
(S, D, R, P, F, H, SC and BC), with 55 descriptors on the original 3,830-recording
Native30 cohort. BC is the fixed center-eight-second median squared
bicoherence statistic at the predeclared 500/750/1,250-Hz coupling coordinates;
it is not a scan selecting the most discriminative frequency triplet.
All five training-group caps and the three distinct grouped evaluation
protocols are retained. There are 103,530 completed model evaluations:
73,950 primary and 29,580 diagnostic cells. The 3,570 structurally ineligible
cells are recorded as omissions, not silently treated as completed trials.

Each fitted fold uses training-derived imputation and scaling, ridge parameter
10, a fixed threshold of 0.5 and uncalibrated linear scores. Primary balanced
accuracy averages correctness within groups, then equally across sources within
each class, then across the two classes; held-source experiments are equally
weighted. AUC is calculated within individual fitted folds, never by pooling
scores from different models. Repeated opposite-class test occurrences do not
constitute additional independent recordings.

\begin{table}[htbp]\centering\small
\begin{tabular}{lrrr}\toprule
Feature set & Generator holdout & Human-source holdout & Grouped descriptive\\\midrule
S & 61.39 & 60.95 & 70.26\\
D & 59.02 & 62.40 & 63.04\\
R & 62.00 & 61.44 & 62.61\\
P & 52.84 & 54.48 & 54.90\\
F & 56.72 & 59.27 & 67.09\\
H & 50.42 & 57.51 & 58.86\\
SC & 52.20 & 60.28 & 63.01\\
BC & 44.62 & 45.96 & 49.62\\
S+D+R+P & 65.79 & 65.65 & 74.30\\
S+D+R+P+BC & 65.05 & 65.19 & 74.20\\
S+D+R+P+F+H+SC & 61.63 & 67.99 & 75.66\\
S+D+R+P+F+H+SC+BC & 61.40 & 67.65 & 75.80\\
\bottomrule\end{tabular}
\caption{Native30 group/source-balanced accuracy (\%) at the all-cap configuration. Protocols have distinct estimands; no cross-column causal comparison is implied.}
\label{tab:bc-allcap}\end{table}
\begin{table}[htbp]\centering\small
\begin{tabular}{lrrrrr}\toprule
Protocol (all eight families) & 25 & 50 & 100 & 200 & All\\\midrule
Generator holdout & 62.47 & 59.30 & 59.41 & 61.02 & 61.40\\
Human-source holdout & 65.43 & 66.52 & 67.07 & 67.69 & 67.65\\
Grouped descriptive & 74.77 & 76.43 & 75.95 & 75.78 & 75.80\\
\bottomrule\end{tabular}
\caption{All-eight-family training-group-cap ablation. Caps are per-source group budgets, not total training-song counts. The complete 3,825-cell table retains all 255 subsets at all five caps.}
\label{tab:bc-caps}\end{table}


Adding BC to the seven-family configuration changes generator-held-out balanced
accuracy from 61.6251\% to 61.3984\% ($-0.2267$ percentage points), and
human-source-held-out balanced accuracy from 67.9920\% to 67.6457\%
($-0.3462$ points). Ordinary grouped descriptive performance changes from
75.6649\% to 75.8033\% ($+0.1384$ points). BC alone reaches 44.6224\%,
45.9584\% and 49.6200\%, respectively. These are descriptive outcomes at the
fixed operating point, not significance tests. Below-chance performance is
reported as observed; labels and score signs are not reversed after testing.
The small grouped-descriptive increase does not establish better transfer.

The report preserves every subset and cap rather than naming a best classifier.
The completed matrix provides no evidence that BC consistently improves
cross-source discrimination. External phase-coupling sensitivity and robust
AI/human attribution remain different claims. In particular, a measurable
coupling statistic need not identify generative provenance.

\paragraph{Admission terminology and evidence limits.}
The legacy reporting schema retains the flag
\texttt{BC\_classifier\_admitted=false}. This must not be read as saying BC was
never supplied to an exploratory fitted model: BC-containing subsets were
explicitly evaluated. It means that the experiment does not establish BC as a
validated general-purpose detector. Report generation performed no fitting,
model selection or threshold tuning. Product hashes and all evaluation receipts
were checked, but the report does not claim a separately implemented numerical
audit of every aggregate metric or confidence intervals over fixed buckets.

\paragraph{Reproducibility.}
The full report is retained locally under
\path{current_results/native30_bc_evaluation_report_v3_20260912/}; its COMMIT
SHA256 is:
\begin{quote}\scriptsize\ttfamily
ad5223a5f1f0fb667fe38a1e0c214b28ab8a8d0b57f01f912521885aa1c0213b
\end{quote}
The underlying evaluation COMMIT is:
\begin{quote}\scriptsize\ttfamily
6d6acecbb11b6fcee72f7ba8730b1a2a26f2797c38d659ca69151074dd3dba06
\end{quote}
The compact CSV in \path{current_results/bc_thesis_tables_v1/} retains all
3,825 primary protocol/subset/cap cells, including denominators and omissions.
The exporter is \path{export_bc_thesis_tables_v1.py}; the source reporter is
\path{summarize_native30_evaluation_bc_v3.py}. Earlier failed report adapters and
earlier thesis versions remain preserved rather than overwritten.
